\section{Related Work}
\label{sec:related_work}

This study connects robotic surface processing, contact-force control, and
model-based reinforcement learning. We position BC-assisted direct TD-MPC2 at their
intersection and examine its operating envelope in force-aware wiping.

\subsection{Robot Wiping, Cleaning, and Polishing}

Surface-processing tasks combine geometric coverage with sustained physical
interaction. Demonstration-driven systems structure this problem through
learned interaction primitives or reusable representations. Probabilistic
Surface Interaction Primitives (ProSIP) conditions compliant task-space
motions on tool-centre-point (TCP) pose, local surface geometry, and
context~\cite{unger2024prosip}; Adaptive Wiping combines few-shot imitation,
force--torque feedback, and pre-trained object representations to adapt
contact-rich execution across task variations~\cite{tsuji2025adaptivewiping}.

Learned contact policies often retain structure in either the action space or
the execution layer. Variable Impedance Control in End-Effector Space (VICES)
treats variable end-effector impedance as a reinforcement-learning (RL)
action space, evaluates it on path following, door opening, and surface wiping,
and reports simulation-to-real transfer
\cite{martinmartin2019vices}. Beltran-Hernandez \emph{et al.} combine RL with
modified parallel position--force control or, separately, admittance control,
and add a real-robot fail-safe
\cite{beltranhernandez2020force}. Zhu \emph{et al.} couple an image-based
policy for equilibrium-point and variable-impedance commands with a faster
momentum-observer and contact-aware null-space controller
\cite{zhu2022contactsafe}. Force-Guided Exploration for Robust Contact-Rich
Manipulation under Uncertainty (FORGE) instead conditions assembly policies on
a maximum allowable force and uses dynamics randomization for transfer from
simulation to hardware~\cite{noseworthy2025forge}. These methods show how
impedance, feedback, and force conditioning can structure learned contact
behaviour. Our comparison places TD-MPC2 and proximal policy optimization
(PPO) on a common direct Cartesian-delta interface, with adaptive admittance as a separate classical
reference.

The closest methods differ from ours along explicit modelling and deployment
axes. VICES uses model-free RL to select variable impedance, whereas our
controller learns latent dynamics and issues Cartesian displacements directly.
Kleff \emph{et al.} incorporate measured force/torque feedback through an explicit
robot-actuation model; our latent model is conditioned on scalar measured
normal force but contains no analytical actuation or contact-force evolution
model~\cite{kleff2022forcefeedbackmpc}. FORGE conditions policies on a maximum
allowable force and uses dynamics randomization for sim-to-real assembly;
our study instead examines first-pass wiping in simulation and audits the
numerical support for each reported outcome~\cite{noseworthy2025forge}. The
distinguishing combination is therefore a learned world model, direct Cartesian
deployment, and claim-level numerical auditing, rather than a new impedance
law, runtime shield, or hardware-transfer method.

Liu \emph{et al.} study quality-critical wiping over randomised curvatures,
frictions, and waypoints using a visual-language-model curriculum and a
checkpoint-bounded reward~\cite{liu2025wipingvlm}. Their 20~Hz policy sends
pose commands to an operational-space controller, with force tracked around a
60~N target. Our complementary setting isolates direct learned Cartesian
control at 5, 8, and 12~N across fixed simulation blocks, enabling trace-level
analysis of completion, force regulation, and transfer.

Curved-surface systems commonly treat tool alignment and force regulation as
a coupled problem. Armesto \emph{et al.} reproduce a wiping policy learned on
planes over an unknown curved surface by estimating the contact constraint and
using force or torque feedback to align the tool~\cite{armesto2018constraintaware}.
Li \emph{et al.} infer pose deviation from contact measurements, lift a
predefined planar path onto an unknown three-dimensional workpiece, and close
a separate normal-force loop~\cite{li2023unknownsurface}. These methods expose
a distinction that matters in our results: our registered frame supplies a
local normal, but the policy action in Eq.~\eqref{eq:task_frame_map} controls
translation only. It cannot directly correct tool posture as the normal turns.

\subsection{Hybrid Position/Force and Impedance Control}

Hybrid position--force control separates motion and force objectives along
different task directions~\cite{raibert1981hybrid}, while impedance control
shapes the dynamic relation between motion and contact
forces~\cite{hogan1985impedance}. Admittance control instead maps interaction
force error into motion of a compliant reference. Impedance and admittance
offer complementary interaction behaviours~\cite{ott2010unified}. We include
an adaptive admittance baseline under the same time-causal observation and
actuator action box. Its force-tier dependence motivates target-wise
compound-pass reporting.

Force-feedback model-predictive control (MPC) provides a closely related
optimization-based alternative. Kleff \emph{et al.} introduce measured
force/torque feedback by augmenting the predictive formulation with robot-actuation
dynamics, and validate the controller on a humanoid upper body
\cite{kleff2022forcefeedbackmpc}. Jordana \emph{et al.} estimate online the
discrepancy between model-predicted and measured force, avoiding an explicit
model of contact-force evolution, and demonstrate force tracking on a
torque-controlled manipulator~\cite{jordana2024forcefeedbackmpc}. Both use
dedicated force-feedback MPC formulations. We address the complementary
question of what a learned actor and world model can achieve through a bounded
Cartesian-displacement interface without a separate force-control layer.

Runtime shields modify policy proposals before
execution~\cite{alshiekh2018shielding}, and control barrier functions provide
a route to forward-invariance conditions
\cite{ames2019cbf,cheng2019rlcbf}. These mechanisms encode safety through
explicit runtime structure. Our direct deployment instead evaluates learned
Cartesian actions without an external shield or force-dependent
post-processing. The resulting 100~Hz force traces provide empirical
exceedance observations for the numerical configuration in which they were
recorded; formal invariance lies outside the study, and
Sec.~\ref{sec:credibility_audit} examines their setting sensitivity.

Contact-rich simulation also depends on contact models, physical relaxations,
and numerical solution methods whose application-level effects need not be
interchangeable~\cite{lelidec2024contact}. This motivates a claim-level
sensitivity study rather than treating a force tail recorded under one
configuration as numerically invariant.

\subsection{Direct Model-Based Reinforcement Learning}

Temporal-difference model-predictive control (TD-MPC) and TD-MPC2 combine
learned latent dynamics, value prediction, and
model-predictive control for continuous
tasks~\cite{hansen2022tdmpc,hansen2024tdmpc2}. TD-MPC2 uses a learned actor to
guide online trajectory optimization through its latent world model. Our
controller retains actor-guided MPPI but removes the earlier force-threshold
router: the actor provides proposals in every regime, while continuous causal
memberships adapt the planning objective and search budget. A same-checkpoint
fixed-planner ablation therefore isolates regime adaptation from the
capability already stored in the actor and world model.
PPO~\cite{schulman2017ppo} provides a matched direct model-free comparator
under a common interface and transition budget.

Alternative designs deliberately retain more structure. Residual reinforcement
learning learns corrections around a conventional controller
\cite{johannink2019residual}; parameterized action primitives and Manipulation
Primitive-augmented Reinforcement Learning (MAPLE) shorten
the decision horizon through learned selection of closed-loop skills
\cite{dalal2021raps,nasiriany2022maple}; and impedance primitive-augmented
hierarchical reinforcement learning (IMP-HRL) combines primitive
selection with variable-stiffness execution~\cite{tahmaz2025imphrl}. Against
these structured alternatives, we characterize the empirical operating
envelope of direct actor/world-model control under controlled variation and
weak-curvature transfer.
